\chapter{Примеры структур данных}
\label{appendix:examples}

В приложении приведены сокращенные примеры структур данных, используемых в MVP. Они демонстрируют связь между чанком, результатом поиска, prompt и ответом с источниками.

\begin{lstlisting}[caption={Пример JSONL-записи чанка}]
{
  "chunk_id": "technical_008::chunk_0004",
  "doc_id": "technical_008",
  "title": "Требования к сети Wi-Fi и IT-интеграции",
  "category": "technical",
  "source_path": "data/raw/02_Технические документы/...",
  "heading_path": ["Требования к Wi-Fi"],
  "text": "Фрагмент документа с требованиями к покрытию сети..."
}
\end{lstlisting}

\begin{lstlisting}[caption={Пример результата dense retrieval}]
{
  "rank": 1,
  "chunk_id": "technical_008::chunk_0004",
  "score": 0.8123,
  "title": "Требования к сети Wi-Fi и IT-интеграции",
  "category": "technical"
}
\end{lstlisting}

\begin{lstlisting}[language=bash,caption={Схема prompt для генерации ответа}]
Вопрос пользователя:
<вопрос>

Контекст из базы знаний:
[Источник 1]
chunk_id: <chunk_id>
title: <title>
text:
<текст фрагмента>

Сформируй ответ по вопросу пользователя.
Используй только предоставленный контекст.
В конце добавь раздел "Источники" со списком chunk_id.
\end{lstlisting}

\begin{lstlisting}[language=bash,caption={Схема ответа с источниками}]
Ответ:
<краткий ответ на русском языке по найденному контексту>

Источники:
- technical_008::chunk_0004
- technical_007::chunk_0002
\end{lstlisting}

Приведенные примеры являются сокращенными и предназначены для пояснения структуры данных. Реальные записи могут содержать дополнительные metadata: диапазоны страниц, листы XLSX, номера строк, block identifiers и служебные поля отчетов.

%%% Local Variables:
%%% TeX-engine: xetex
%%% End:
